\subsection{Постановка задачи}

Цель курсовой работы – проектирование и реализация файлового
менеджера с возможностью просмотра файлов на устройствах в
одной локальной сети в виде кроссплатформенного настольного
приложения.

\begin{itemize}
    \item модуль утилит;
    \item модуль настроек;
    \item модуль файловой системы;
    \item модуль сетевых операций;
    \item модуль интерфейса пользователя;
\end{itemize}

При разработке программного средства будет использована
интегрированная среда разработки CLion \cite{clionDOC}
с использованием языка программирования C++ \cite{cppDOC}
и платформы MinGW \cite{mingwDOC} для сборки проекта.

Проект будет разработан с использованием модульной архитектуры \cite{modularArch},
где каждый модуль будет отвечать за свою функциональность:

\begin{itemize}
    \item Разделение ответственности -- каждый модуль имеет чётко определённую область ответственности;
    \item Инкапсуляция -- модули скрывают детали реализации и предоставляют интерфейсы для взаимодействия;
    \item Масштабируемость -- модульная структура позволяет легко добавлять новую функциональность;
    \item Поддерживаемость -- изолированные модули упрощают отладку и модификацию кода;
\end{itemize}